\SourceLeaf{44}
\fancyhead[L]{\small METHODUS PECULIARIS ELEVATIONEM POLI DETERMINANDI}
\renewcommand{\Article}[1]{\par\medskip\Needspace{5\baselineskip}\pdfbookmark[1]{Art. #1}{poli1808-art-#1}\begin{center}\textbf{#1.}\end{center}\nobreak}
\pdfbookmark[0]{Methodus peculiaris elpvationem poli determinandi}{poli1808-body}
\begin{center}\large METHODUS PECULIARIS\par\medskip
ELPVATIONEM POLI DETERMINANDI.\par\medskip\rule{.18\linewidth}{.3pt}\end{center}
\Article{1}
Methodos usitatissimas, quas observatores itinerantes ad determinationes geographicas adhibere solent, observationibus Solaribus inniti constat, altitudinibus correspondentibus determinationi temporis, altitudinibusque meridianis vel circummeridianis determinationi elevationis poli inservientibus. Pluribus utique nominibus hocce observationum genus se commendat, iis praesertim, qui ad altitudines mensurandas sextante Hadleyano utuntur: multo enim facilius est altitudines Solis hocce instrumento accurate observare, quam altitudines stellarum, limbique divisiones nocte minus commode dignoscuntur, quam interdiu; calculus praeterea altitudinum correspondentium et circummeridianarum tam simplex est, ut a quolibet etiam in mathesi parum versato facile addisci possit. --- Quodsi forte propter nubes aliave impedimenta altitudines correspondentes assequi non licuerit, observationes culminationi proximae cum aliis magis inde remotis combinari solent, horologiique a tempore vero deviatio, nec non altitudo poli plerumque per methodos indirectas, e. g. per eam quae a \textsc{Cornelio Douwes} nomen accepit, inde derivari.

Quamquam vero nullum sit dubium, quin observationes Solares nocturnis praeferendae sint, quoties quidem illarum compotem fieri licet: tamen posteriores haud omnino negligi debere censemus. Saepissime coelum, postquam interdiu nubibus tectum fuit, ingruente nocte defaecatur, aut saltem stellas subinde
\SourceLeaf{45}conspicere concedit: saepissime etiam viator, cui interdiu itineri prosequendo intento observationibus vacare non licuerat, per determinationem geographicam loci, in quo pernoctat, bene meritum se praestare posset, si modo in observandis stellis dexteritatem sufficientem sibi comparavisset. Adnumeremus adhuc hisce rationibus incertitudinem observationum Solarium in mensibus hibernis, ubi nostris quidem regionibus Sol parum ultra vapores horizontis exsurgit, porro frequentiores observationum perturbationes a spectatoribus importunis, concussiones horizontis artificialis (praesertim si mercuriali utaris) a curribus praetervectis, a quibus per silentium noctis magis tutus eris, faterique oportebit, observationes stellarum perdignas esse, quae diligentius excolantur, atque in usum vocentur.

Quantam utilitatem hocce observationum genus \textit{in itineribus maritimis} praestare possit, facile patebit, si perpendatur, quanti momenti sit, ut navis locus vel quotidie si fieri possit exploretur, a qua cognitione haud raro hominum salus et vita pendet. Haud equidem diffitendum est, in mari quoque observationes nocturnas difficultatibus peculiaribus premi, quum horizon maris, a quo altitudines capere oportet, nocte minus commode distinguatur: sed constat quoque, huic incommodo remedium afferri posse, si sextans tubo optico aperturae maioris muniatur.
\Article{2}
Quae in art. praec. in laudem observationum stellarum allata sunt, vim adhuc maiorem nanciscuntur, quum hoc modo brevissimo temporis intervallo, puta per observationes intra pauca minuta prima instituendas, non modo altitudo poli, sed insuper tempus mira praecisione determinari possit. Ad hunc finem duas tantummodo altitudines stellarum diversarum quarumcunque (quarum tamen positiones notae supponuntur) observare, horologiique tempora respondentia notare oportet: motus horologii diurnus ea tantummodo praecisione notus esse debet, ut intervallum temporis satis exacte in tempus sidereum converti possit. Certo ni fallimur hoc problema inter utilissima astronomiae nauticae referendum est, adeoque satis est mirandum, quod ab iis qui supra hac re scripserunt penitus neglectum esse videtur. Clar. \textsc{Kraft} in \textit{Actis novis Acad. Petropol.} T. XIII. ex altitudinibus duarum stellarum elevationem poli determinare docuit: sed, quanam ratione inductus nescio, problema generale per conditionem specialem limitavit, supponendo scilicet, ambas observationes eodem temporis momento factas
\SourceLeaf{46}esse. Tunc itaque requiruntur duo observatores, duo instrumenta, exacta observationum congruentia, quam sane haud ita facile assequi licebit. Verum enim vero haecce conditio prorsus inutilis et superflua est. Certo enim navis bene instructa carere nequit horologio, quod intervallum paucorum minutorum exacte dimetiri valeat; neque etiam motus navis intra tantillum tempus sensibilem positionis geographicae mutationem producere poterit; denique calculus problemati solvendo inserviens per conditionem illam nihilo fit simplicior, uti statim monstrabimus.
\Article{3}
Nullo negotio perspicietur, problema in forma ea, in qua a clar. \textsc{Kraft} tractatum est, ita quoque enunciari posse. `Ex data positione duorum punctorum in superficie sphaerica respectu alicuius circuli maximi (qui hic est aequator) invenire positionem puncti tertii, cuius distantiae ab illis punctis sunt datae.' Puncta scilicet priora dabuntur per positiones duarum stellarum observatarum in sphaera coelesti, puta per harum ascensiones rectas et declinationes: punctum tertium autem respondebit ipsi zenith loci observationis, ipsiusque declinatio elevationem poli, ascensio recta punctum culminans aequatoris, adeoque tempus sidereum definiet. Solutio huius problematis statim obvia per tria triangula sphaerica perficitur. In primo, quod formabitur inter duo puncta data atque polum aequatoris, data sunt duo latera cum angulo incluso, unde derivabuntur latus tertium, atque alteruter duorum angulorum reliquorum, e. g. is qui puncto primo adiacet. In triangulo secundo, quod inter duo puncta data punctumque tertium quaesitum formabitur, data iam sunt latera omnia, unde derivabitur unus ex angulis, is scilicet qui itidem puncto primo adiacet, siquidem angulum eidem adiacentem in triangulo primo adoptaveramus. Summa vel differentia horum duorum angulorum constituet angulum notum in triangulo tertio, inter punctum primum, tertium polumque aequatoris formato, in quo insuper duo latera adiacentia cognita sunt: hinc denique angulus ad polum, punctique tertii distantia a polo (quae est complementum elevationis poli) derivabuntur. Haecce methodus iam ab astronomis saeculi XVI in usum vocata est: vid. e. g. \textit{Tychonis Astronomiae instauratae progymnasmata} p. 211 sqq., ubi permultarum stellarum (interque eas etiam novae in Cassiopea) positiones e distantiis a binis stellis positione datis determinantur.

\SourceLeaf{47}Simul hinc manifestum est, stricte loquendo problema duas semper solutiones admittere: si enim in superficie sphaerica e duobus punctis datis tamquam centris descripti supponantur duo circuli radiis, qui resp. distantiis propositis aequales sunt, hi circuli generaliter loquendo duobus punctis se secabunt, punctumque tertium quaesitum in utraque intersectione concipi poterit. Nihilominus in praxi ambiguitas hinc oriri nequit. Scilicet si circulus maximus per punctum primum et secundum ductus concipitur, superficiem sphaericam in duo hemisphaeria dirimens, dubium esse non potest, utrum punctum tertium (ipsum zenith) polusque arcticus in eodem hemisphaerio iacere debeant, an in oppositis: patet enim, casum primum locum habere debere, si excessus ascensionis rectae stellae eius, quae in observatione fuit ad laevam (respectu stellae alterius) supra ascensionem rectam stellae alterius, quae respectu prioris fuit ad dextram (adiectis si opus fuerit $360$ gradibus) fuerit inter $0$ et $180^0$; contra, casum posteriorem valere. si excessus iste fuerit inter $180^0$ et $360^0$.

Ceterum huic ipsi solutionis ambiguitati, problemati a priori inhaerenti, attribui debet calculi prolixitas, qua solutio directa laborat; in sequentibus tamen illam notabiliter contrahere docebimus. Quodsi autem solutionem indirectam adhibere non recusemus, satis expedita adstrui potest, cuius explicationem ad aliam occasionem nobis reservamus.

Hactenus de casu speciali, ubi duae observationes simultaneae sunt, diximus: problema generale ad hunc casum sponte reducitur sequenti modo. Concipiatur pro stella, quae secundo loco observata est, alia fictitia eandem cum illa declinationem habens, sed ascensionem rectam tanto minorem, quanto temporis siderei intervallo observatio secunda primam sequuta est. Manifesto haec stella fictitia tempore observationis primae eandem altitudinem attigisset, quam stella revera observata habuit tempore observationis secundae, unde illa huic substituta iam calculus observationibus simultaneis superstruendus erit.

Haec omnia e considerationibus obviis quidem sed geometricis petita sunt: quamobrem speramus, solutionem directam solis operationibus analyticis innixam pluribus haud ingratam fore. Nova certe hinc prodibit confirmatio, quod vix quidquam e considerationibus geometricis hauriatur, quod non possit aeque concinne per analysin erui, si modo scite tractetur.
\SourceLeaf{48}\Article{4}
Sit $\varphi$ elevatio poli, $\alpha$ et $\alpha'$ ascensiones rectae, $\delta$ et $\delta'$ declinationes duarum stellarum, $\gamma$ et $\gamma'$ ascensiones rectae punctorum culminantium aequatoris resp. tempore observationis primae et secundae, sive quod idem est tempora ipsa siderea in gradus conversa, denique $h$ altitudo observata stellae primae, $h'$ altitudo stellae secundae. Manifesto $\gamma-\alpha$, $\gamma'-\alpha'$ erunt anguli horarii, duabus observationibus respondentes: statuemus $\gamma-\alpha=\lambda$, $\gamma'-\alpha'=\lambda-\theta$, unde $\theta=\alpha'-\alpha-(\gamma'-\gamma)$ erit quantitas cognita, quoniam $\gamma'-\gamma$ aequalis est intervallo temporis inter duas observationes in tempus sidereum ac dein in gradus converso. Problematis itaque solutio hisce duabus aequationibus superstruenda erit:
\[
\begin{array}{rl}
{}[1]\quad&\sin h=\sin\delta\sin\varphi+\cos\delta\cos\varphi\cos\lambda\\
{}[2]\quad&\sin h'=\sin\delta'\sin\varphi+\cos\delta'\cos\varphi\cos(\lambda-\theta)
\end{array}
\]
Quodsi iam ex his aequationibus vel incognitam $\lambda$ vel incognitam $\varphi$ eliminare susciperemus, ad aequationem nimis complicatam deferremur: priorem quidem viam sequutus est clar. \textsc{Kraft} in diss. supra citata, sed nostro quidem iudicio solutio, quam hoc modo eruit, longior adhuc molestiorque est ea, quae immediate e consideratione trium triangulorum derivatur, etiamsi illa ad solam determinationem ipsius $\varphi$ limitetur, determinatione temporis neglecta. Praestabit itaque, incognitam aliquam novam introducere, eiusque adiumento utramque $\varphi,\lambda$ eliminare. Ad quam apte eligendam sequens observatio viam nobis sternet. Fit
\[
\begin{aligned}
&(\sin\delta\sin\varphi+\cos\delta\cos\varphi\cos\lambda)^2+(\cos\delta\sin\varphi-\sin\delta\cos\varphi\cos\lambda)^2\\
&\qquad=\sin\varphi^2+\cos\varphi^2\cos\lambda^2=1-\cos\varphi^2\sin\lambda^2
\end{aligned}
\]
adeoque per aequationem $1$
\[
(\cos\delta\sin\varphi-\sin\delta\cos\varphi\cos\lambda)^2+\cos\varphi^2\sin\lambda^2=\cos h^2,
\]
sive
\[
\left(\frac{\cos\delta\sin\varphi-\sin\delta\cos\varphi\cos\lambda}{\cos h}\right)^2+\left(\frac{\cos\varphi\sin\lambda}{\cos h}\right)^2=1
\]

Statuere itaque licebit
\[
\begin{array}{rl}
{}[3]\quad&\dfrac{\cos\delta\sin\varphi-\sin\delta\cos\varphi\cos\lambda}{\cos h}=\cos u\\[4pt]
{}[4]\quad&\dfrac{\cos\varphi\sin\lambda}{\cos h}=\sin u
\end{array}
\]

\SourceLeaf{49}Iam si aequatio $3$ cum $1$ combinatur, eruimus
\[
\begin{array}{rl}
{}[5]\quad&\sin\varphi=\sin\delta\sin h+\cos\delta\cos h\cos u\\
{}[6]\quad&\cos\varphi\cos\lambda=\cos\delta\sin h-\sin\delta\cos h\cos u
\end{array}
\]

Tribuamus aequationi $2$ formam hancce
\[
\sin h'=\sin\delta'\sin\varphi+\cos\theta\cos\delta'\cos\varphi\cos\lambda+\sin\theta\cos\delta'\cos\varphi\sin\lambda
\]
substituamusque pro $\sin\varphi$, $\cos\varphi\cos\lambda$, $\cos\varphi\sin\lambda$ valores suos ex aequationibus $5,\ 6,\ 4$, unde fiet
\[
\begin{aligned}
&\sin h'-\sin h\sin\delta\sin\delta'-\sin h\cos\theta\cos\delta\cos\delta'-\cos u\cos h\cos\delta\sin\delta'\\
&\quad+\cos u\cos h\cos\theta\sin\delta\cos\delta'-\sin u\cos h\sin\theta\cos\delta'=0
\end{aligned}
\]

Haec aequatio, statuendo
\[
{}[7]\qquad\frac{\cos\delta\sin\delta'-\cos\theta\sin\delta\cos\delta'}{\sin\theta\cos\delta'}=\operatorname{cotang}v
\]
transit in sequentem
\[
\begin{aligned}
&\sin h'-\sin h\sin\delta\sin\delta'-\sin h\cos\theta\cos\delta\cos\delta'\\
&\qquad-\cos h\sin\theta\cos\delta'(\cos u\operatorname{cotang}v+\sin u)=0
\end{aligned}
\]

Quodsi itaque statuimus $v-u=w$, habebimus
\[
{}[8]\qquad\cos w=\frac{\sin v(\sin h'-\sin h\sin\delta\sin\delta'-\sin h\cos\theta\cos\delta\cos\delta')}{\cos h\sin\theta\cos\delta'}
\]
unde determinabitur $w$, atque hinc $u=v-w$.

Postquam angulus $u$ inventus est, deducetur $\lambda$ e combinatione aequationum $4$ et $6$, unde oritur
\[
{}[9]\qquad\operatorname{tang}\lambda=\frac{\cos h\sin u}{\cos\delta\sin h-\sin\delta\cos h\cos u}
\]

Ex $\lambda$ habebitur $\gamma=\alpha+\lambda$, quo angulo in tempus converso status horologii innotescet. Denique per combinationem aequationum $4$ et $5$ prodit
\[
{}[10]\qquad\operatorname{tang}\varphi=\frac{\sin\lambda\,(\sin\delta\sin h+\cos\delta\cos h\cos u)}{\cos h\sin u}
\]
cui formulae, si placet, calculi confirmandi caussa ipsa aequatio $4$ adiungi poterit.
\SourceLeaf{50}\Article{5}
Circa solutionem in art. praec. traditam quasdam adhuc observationes adiicimus.

Quum e formula $8$ pro angulo $w$, per cosinum suum determinando, duo valores prodeant, alter positivus, alter negativus, duas inde solutiones diversas emergere patet, ut iam supra monuimus: attamen utra sit vera, facile sequenti modo dignoscitur. Ostendi potest, sinum differentiae inter azimuthum circuli verticalis, in quo stella prima observata est, atque azimuthum circuli verticalis, in quo observata est stella secunda, sive exactius, sinum excessus azimuthi prioris supra posterius a laeva ad dextram mensurati, fieri
\[
=\frac{\sin\theta\cos\delta'\sin w}{\cos h'\sin v}
\]
Iam quum $\cos\delta',\ \cos h'$ natura sua sint quantitates positivae, patet, $\sin w$ vel eodem signo affectum esse debere, quod habet $\frac{\sin\theta}{\sin v}$, vel opposito, prout circulus verticalis prior posteriori vel ad dextram fuerit vel ad laevam, de qua re incertitudo adesse nequit, quum observationes in circulis verticalibus vel prope conspirantibus, vel prope oppositis factae, ad praxin omnino non sint idoneae, adeoque sedulo evitari debeant, uti infra fusius exponemus.

Ceterum solutio adhuc aliquantulum contrahi calculoque numerico magis accommodari potest per introductionem duorum angulorum auxiliarium. Determinando scilicet angulum $F$ per aequationem
\[
{}[11]\qquad\operatorname{tang}F=\frac{\operatorname{tang}\delta'}{\cos\theta}
\]
aequationes $7$ et $8$ transeunt in has:
\[
\begin{array}{rl}
{}[12]\quad&\operatorname{tang}v=\dfrac{\cos F\operatorname{tang}\theta}{\sin(F-\delta)}\\[5pt]
{}[13]\quad&\cos w=\dfrac{\cos v\operatorname{tang}h}{\operatorname{tang}(F-\delta)}\left(\dfrac{\sin h'\sin F}{\sin h\sin\delta'\cos(F-\delta)}-1\right)
\end{array}
\]

Perinde introducendo angulum auxiliarem $G$ talem, ut sit
\[
{}[14]\qquad\operatorname{tang}G=\frac{\operatorname{tang}h}{\cos u}
\]
aequationes $9$ et $10$ in has transibunt:
\[
\begin{array}{rl}
{}[15]\quad&\operatorname{tang}\lambda=\dfrac{\cos G\operatorname{tang}u}{\sin(G-\delta)}\\[5pt]
{}[16]\quad&\operatorname{tang}\varphi=\cos\lambda\operatorname{cotang}(G-\delta)
\end{array}
\]
\SourceLeaf{51}\Article{6}
Astronomi practici saepius queruntur, methodos a theoreticis propositas non semper ad usus practicos idoneas esse, vel saltem non praestare, quae inde exspectaverint. Quo melius talem suspicionem a methodo nostra avertamus simulque praeceptorum praecedentium usum illustremus, exemplum aliquod ab observationibus non fictis desumtum fusius exsequi conveniet. Sextante Troughtoniano, cuius radius decem pollicum, observavi in specula nostra Aug. $21$ huius anni altitudinem duplicem non correctam stellae $\alpha$ Aquilae $91^0\ 34'\ 10''$, quum horologium Sheltonianum monstraret $20^{\mathrm h}\ 40^{\mathrm m}\ 8^{\mathrm s}$; altitudinem aequalem assequuta est stella $\alpha$ Andromedae tempore horologii $20^{\mathrm h}\ 46^{\mathrm m}\ 59^{\mathrm s}$; stella prior observata est in circulo verticali versus dextram, posterior versus laevam; motus horologii diurnus cum tempore sidereo vero exacte congruit. Purgato angulo observato ab errore indicis prodit angulus verus $91^0\ 31'\ 36''$, adeoque altitudo apparens $45^0\ 45'\ 48''$, hinc altitudo vera $45^0\ 44'\ 52''6=h=h'$. Stellarum positiones, aberratione et nutatione rite applicatis, inventae sunt hae:
\[
\begin{array}{rcl@{\quad}rcl}
\alpha&=&295^0\ 22'\ 6''6&\delta&=&+8^0\ 22'\ 43''1\\
\alpha'&=&359\ 38\ 18.5&\delta'&=&+28\ 2\ 13.4
\end{array}
\]
Differentia temporum $6'\ 51''$ in arcum conversa producit $\gamma'-\gamma=1^0\ 42'\ 45''$, unde $\theta=62^0\ 33'\ 26''9$. Calculus ipse dein ita se habet (factor $\frac{\sin h'}{\sin h}$ in formula $13$ propter altitudinum aequalitatem manifesto hic excidit):
\[
\begin{array}{lrl}
\log\operatorname{tang}\delta'\quad .\ .\ .\ .&9.7263516&\\
\log\cos\theta\quad .\ .\ .\ .\ .&9.6635677&\\\cline{1-2}
\log\operatorname{tang}F\quad .\ .\ .\ .&0.0627839,&\text{unde }F=49^0\ 7'\ 37''70\\
&&\text{atque }F-\delta=40^0\ 44'\ 54''60
\end{array}
\]
\[
\begin{array}{lrl}
\log\operatorname{tang}\theta.\quad .\ .\ .\ .&0.2845877&\\
\log\cos F\quad .\ .\ .\ .\ .&9.8158318&\\
\text{Comp. }\log\sin(F-\delta)\quad .&0.1802598&\\\cline{1-2}
\log\operatorname{tang}v\quad .\ .\ .\ .\ .&0.2856793&\text{unde }v=62^0\ 36'\ 58''79
\end{array}
\]
\[
\begin{array}{lrl}
\log\sin F\quad .\ .\ .\ .\ .&9.8786157&\\
\text{Comp. }\log\sin\delta'\quad .\ .\ .&0.3278629&\\
\text{Comp. }\log\cos(F-\delta)\quad .&0.1205703&\\\cline{1-2}
&0.3270489&=\log 2.123483=\log n
\end{array}
\]
